\begin{problem}
    已知方阵$\mat{A}$的一个特征值为$10$，则$2\mat{A}^3+\mat{A}^2-77\mat{E}$必有一个特征值是\ans{$2023$}.
\end{problem}
\begin{solution}*
    $2\mat{A}^3+\mat{A}^2-77\mat{E}$对应的特征值为$2\lambda^3+\lambda^2-77=2\times 10^3+10^2-77=2023$.
\end{solution}

\begin{problem}
    已知二次型$f(x_1,x_2,x_3)=2x_1x_2+2x_1x_3+2x_2x_3$，则其正惯性指数为\ans{$1$}.
\end{problem}
\begin{solution}
    写出二次型的矩阵
    $\mat{A}$，求$\mat{A}$的所有特征值.
    \[
    \mat{A}=\begin{pmatrix}
        0 & 1 & 1\\
        1 & 0 & 1\\
        1 & 1 & 0
    \end{pmatrix},\ 
    \abs{\mat{A}-\lambda\mat{E}}=\begin{vmatrix}
        -\lambda & 1 & 1\\
        1 & -\lambda & 1\\
        1 & 1 & -\lambda
    \end{vmatrix}=-(\lambda+1)^2(2-\lambda)
    \]

    特征值为$\lambda_1=\lambda_2=-1<0,\ \lambda_3=2>0$. 所以正惯性指数为$1$.
\end{solution}

\begin{problem}
    已知方阵$\mat{A}$能够相似对角化，且其五个特征值分别为：$3,\ 4,\ 5,\ 5,\ 5$，则$\rank{\mat{A}-5\mat{E}}=$\ans{$2$}.
\end{problem}
\begin{solution}
\begin{itemize}
    \item $5$是一个三重根，若要使$\mat{A}$能够相似对角化，则$(\mat{A}-5\mat{E})\mat{X}=0$基础解系所含向量必须有$3$个.
    \item 基础解系所含向量个数即为方程自由变量的个数，方程系数矩阵的秩$\rank{\mat{A}-5\mat{E}}$为非自由变量的个数.
    \item 自由变量个数与非自由变量的个数之和为方程组未知数的个数，即$3+\rank{\mat{A}-5\mat{E}}=5$，所以$\rank{\mat{A}-5\mat{E}}=2$.
\end{itemize}
\end{solution}

\begin{problem}
    已知向量$\vv{\alpha}_1=(1,1,1)^\T$，则与向量$\alpha_1$正交的所有向量所构成向量组的秩为\ans{$2$}.
\end{problem}
\begin{solution}
\begin{itemize}
    \item 设与$\vv{\alpha}_1$正交的向量为$\vv{x}=(x_1,x_2,x_3)^\T$，所以$[\vv{\alpha},\vv{x}]=0$，即$x_1+x_2+x_3=0$.
    \item 解方程可设$x_2,\ x_3$为自由变量$(1,0)^\T,\ (0,1)^\T$，则方程组的基础解系为$(1,1,0)^\T,\ (1,0,1)^\T$.
    \item 基础解系所含向量个数即为向量组的秩，为$2$.
\end{itemize}
\end{solution}

\begin{problem}
    若$\vv{\alpha}_1=(1,2,-1)^\T,\ \alpha_2=(-1,3,2)^\T,\ \vv{\alpha}_3=(4,-2,\lambda)^\T$是向量空间$\mathbb{R}^3$中的一组基，则$\lambda$的取值应满足\ans{$\lambda\neq -6$}.
\end{problem}
\begin{solution}
    基满足$\vv{\alpha}_1,\ \vv{\alpha}_2,\ \vv{\alpha}_3$线性无关，则$\vv{\alpha}_1,\ \vv{\alpha}_2,\ \vv{\alpha}_3$组成的矩阵的行列式不为$0$，即
    \[\abs{\vv{\alpha}_1,\ \vv{\alpha}_2,\ \vv{\alpha}_3}\neq 0,\ 
    \begin{vmatrix}
        1 & -1 & 4\\
        2 & 3 & -2\\
        -1 & 2 & \lambda
    \end{vmatrix}\xlongequal[r_3+r_1]{r_2-2r_1}
    \begin{vmatrix}
        1 & -1 & 4\\
        0 & 5 & 10\\
        0 & 1 & \lambda+4
    \end{vmatrix}=5\lambda+20+10\neq 0\]

所以当$\lambda\neq -6$时，$\vv{\alpha}_1,\ \vv{\alpha}_2,\ \vv{\alpha}_3$是$\mathbb{R}^3$的一组基.
\end{solution}

\begin{problem}
    已知矩阵$\mat{A}=\begin{pmatrix}
        -2 & 1 & 1\\
        0 & x & 0\\
        -4 & 1 & 3
    \end{pmatrix}$与
    $\mat{B}=\begin{pmatrix}
        y & 0 & 0\\
        0 & 2 & 0\\
        0 & 0 & 2
    \end{pmatrix}$相似，则$x+y=$\ans{$1$}.
\end{problem}
\begin{solution}
    $\mat{A}$与$\mat{B}$相似，则$\begin{cases}
        \text{tr}(\mat{A})=\text{tr}(\mat{B})\\
        \abs{\mat{A}}=\abs{\mat{B}}
    \end{cases}$，所以有
    $
    -2+x+3=y+2+2,\ \abs{\mat{A}}=\abs{\mat{B}}
    $，解得$x=2,\ y=-1,\ x+y=1$.
\end{solution}

\begin{choice}{B C}
    设$\vv{\alpha}_1,\ \vv{\alpha}_2,\ \vv{\alpha}_3$线性无关，则下列向量组也线性无关的是
\begin{tasks}(2)
    \task $\vv{\alpha}_1-\vv{\alpha}_2,\ \vv{\alpha}_2-\vv{\alpha}_3,\ \vv{\alpha}_3-\vv{\alpha}_1$
    \task $\vv{\alpha}_1+\vv{\alpha}_2,\ \vv{\alpha}_2+\vv{\alpha}_3,\ \vv{\alpha}_3+\vv{\alpha}_1$
    \task $\vv{\alpha}_1-\vv{\alpha}_2-\vv{\alpha}_3,\ \vv{\alpha}_2,\ \vv{\alpha}_3$
    \task $2\vv{\alpha}_1+\vv{\alpha}_2,\ \vv{\alpha}_2+2\vv{\alpha}_3,\ \vv{\alpha}_1+\vv{\alpha}_2+\vv{\alpha}_3$
\end{tasks}
\end{choice}
\begin{solution}
\begin{tasks}
    \task $(\vv{\alpha}_1-\vv{\alpha}_2,\ \vv{\alpha}_2-\vv{\alpha}_3,\ \vv{\alpha}_3-\vv{\alpha}_1)=(\vv{\alpha_1},\vv{\alpha}_2,\vv{\alpha}_3)\begin{pmatrix}
        1 & 0 & -1\\
        -1 & 1 & 0\\
        0 & -1 & 1
    \end{pmatrix},\ \begin{vmatrix}
        1 & 0 & -1\\
        -1 & 1 & 0\\
        0 & -1 & 1
    \end{vmatrix}=0$，所以相关.
    \task $(\vv{\alpha}_1+\vv{\alpha}_2,\ \vv{\alpha}_2+\vv{\alpha}_3,\ \vv{\alpha}_3+\vv{\alpha}_1)=(\vv{\alpha_1},\vv{\alpha}_2,\vv{\alpha}_3)\begin{pmatrix}
        1 & 0 & 1\\
        1 & 1 & 0\\
        0 & 1 & 1
    \end{pmatrix},\ \begin{vmatrix}
        1 & 0 & 1\\
        1 & 1 & 0\\
        0 & 1 & 1
    \end{vmatrix}=2\neq 0$，所以无关.
    \task $(\vv{\alpha}_1-\vv{\alpha}_2-\vv{\alpha}_3,\ \vv{\alpha}_2,\ \vv{\alpha}_3)=(\vv{\alpha_1},\vv{\alpha}_2,\vv{\alpha}_3)\begin{pmatrix}
        1 & 0 & 0\\
        -1 & 1 & 0\\
        -1 & 0 & 1
    \end{pmatrix},\ \begin{vmatrix}
        1 & 0 & 0\\
        -1 & 1 & 0\\
        -1 & 0 & 1
    \end{vmatrix}=1\neq 0$，所以无关.
    \task $2\vv{\alpha}_1+\vv{\alpha}_2,\ \vv{\alpha}_2+2\vv{\alpha}_3,\ \vv{\alpha}_1+\vv{\alpha}_2+\vv{\alpha}_3)=(\vv{\alpha_1},\vv{\alpha}_2,\vv{\alpha}_3)\begin{pmatrix}
        2 & 0 & 1\\
        1 & 1 & 1\\
        0 & 2 & 1
    \end{pmatrix},\ \begin{vmatrix}
        2 & 0 & 1\\
        1 & 1 & 1\\
        0 & 2 & 1
    \end{vmatrix}=0$，所以相关.
\end{tasks}
\end{solution}

\begin{choice}{C}
    下列说法不正确的是
\begin{tasks}(2)
    \task 可逆矩阵一定是满秩矩阵
    \task 等价的矩阵一定有相同的秩
    \task 对角矩阵一定是可逆矩阵
    \task 初等矩阵一定是可逆矩阵
\end{tasks}
\end{choice}
\begin{solution}
\begin{tasks}(2)
    \task $\text{可逆矩阵}=\text{行满秩阵}=\text{列满秩阵}=\text{满秩阵}$.
    \task 等价必等秩.
    \task! 对角阵主对角线元素之和为$0$，则行列式为$0$，不可逆.
    \task! 所有初等矩阵都可逆，因其为单位阵变换行或列近一次得到.
\end{tasks}
\end{solution}

\begin{choice}{A}
    下列哪个矩阵的行列式可能等于零
\begin{tasks}(4)
    \task 对称矩阵
    \task 可逆矩阵
    \task 正交矩阵
    \task 正定矩阵
\end{tasks}
\end{choice}
\begin{solution}
\begin{tasks}(2)
    \task 对称矩阵举例：$\begin{pmatrix}
        0 & 1 & 0\\1 & 0 & 0\\0 & 0 & 0
    \end{pmatrix}$行列式为$0$.
    \task 可逆矩阵行列式一定不为$0$.
    \task 正交矩阵一定是可逆矩阵.
    \task 正交矩阵的行列式为$\pm 1$.
\end{tasks}
\end{solution}

\begin{choice}{D}
    向量$\vv{\beta}=(2,7,2)^\T$在基$\vv{\alpha}_1=(1,2,-1)^\T,\ \vv{\alpha}_2=(-1,1,1)^\T,\ \vv{\alpha}_3=(4,-1,0)^\T$下的坐标为
\begin{tasks}(4)
    \task $(4,2,1)^\T$
    \task $(1,2,4)^\T$
    \task $(2,1,4)^\T$
    \task $(2,4,1)^\T$
\end{tasks}
\end{choice}
\begin{solution}
    设坐标为$x_1,\ x_2,\ x_3$，$\beta=x_1\vv{\alpha}_1+x_2\vv{\alpha}_2+x_3\vv{\alpha}_3$，即
    \[(1,2,-1)^\T x_1+(-1,1,1)^\T x_2+(4,-1,0)^\T x_3=(2,7,2)^\T\]

    即解方程$\mat{A}\vv{x}=\vv{\beta}$，其中$\mat{A}=\begin{pmatrix}
        1 & -1 & 4\\
        2 & 1 & -1\\
        -1 & 1 & 0
    \end{pmatrix}$. 对增广矩阵行变换
    $\begin{pmatrix}
        1 & -1 & 4 & 2\\
        2 & 1 & -1 & 7\\
        -1 & 1 & 0 & 2
    \end{pmatrix}\sim
    \begin{pmatrix}
        1 & 0 & 0 & 2\\
        0 & 1 & 0 & 4\\
        0 & 0 & 1 & 1
    \end{pmatrix}$.
    
    所以坐标为$(2,4,1)^\T$.
\end{solution}

\includepdf[pages={1-12}]{xdwritten}